% =============================================================================
% Bloque 5: Hipótesis y Objetivos (slide 11)
% =============================================================================

\section{Hipótesis y objetivos}

% --- Slide 11: H1 + sub-hipótesis + objetivos ---
\begin{frame}{Hipótesis y objetivos}
\setbeamercolor{block title}{fg=white,bg=uaiorange}
\begin{block}{H1 — Hipótesis principal}
\small
El filtrado geométrico de muestras LLM en el espacio de embeddings produce
mejoras \textbf{estadísticamente significativas} en macro F1 respecto a SMOTE,
en clasificación con pocos ejemplos por clase.
\end{block}

\vspace{0.4em}
\begin{columns}[T]

\begin{column}{0.5\textwidth}
\setbeamercolor{block title}{fg=white,bg=uaiblue}
\begin{block}{Sub-hipótesis verificables}
\scriptsize
\begin{description}
  \item[H1a] Filtrado por distancia $>$ sin filtrar.
  \item[H1b] Ponderación suave $\geq$ filtrado binario.
  \item[H1c] Generaliza a NER sin modificación.
  \item[H1d] Lineales $>$ ensambles de árboles.
\end{description}
\end{block}
\end{column}

\begin{column}{0.5\textwidth}
\setbeamercolor{block title}{fg=white,bg=uaibluedark}
\begin{block}{Objetivos específicos}
\scriptsize
\begin{enumerate}
  \item Diseñar e implementar un sistema de filtros con ponderación suave.
  \item Comparar contra \textbf{9 métodos} (clásicos + modernos) sobre múltiples datasets.
  \item Demostrar generalización a \textbf{NER}.
\end{enumerate}
\end{block}
\end{column}

\end{columns}

\vspace{0.4em}
{\scriptsize\itshape Criterio de aceptación:}
{\scriptsize $\Delta>0$,\quad $p_{\text{Bonf}}<0.05$,\quad win-rate $>50\%$.}
\end{frame}
